\documentclass[letterpaper]{article}
\usepackage[submission]{aaai2027}
\usepackage{times}
\usepackage{helvet}
\usepackage{courier}
\usepackage[hyphens]{url}
\usepackage{graphicx}
\usepackage{natbib}
\urlstyle{rm}
\def\UrlFont{\rm}
\usepackage{caption}
\frenchspacing
\setlength{\pdfpagewidth}{8.5in}
\setlength{\pdfpageheight}{11in}
\usepackage{amsmath}
\usepackage{amssymb}
\usepackage{amsthm}
\usepackage{booktabs}
\usepackage{multirow}
\usepackage{tikz}
\usepackage{placeins}
\emergencystretch=1.2em
\newtheorem{theorem}{Theorem}
\newtheorem{proposition}{Proposition}
\newtheorem{definition}{Definition}
\newtheorem{lemma}{Lemma}
\newtheorem{corollary}{Corollary}
\renewcommand{\qedsymbol}{}
\pdfinfo{
/TemplateVersion (2026.1)
}
\setcounter{secnumdepth}{2}

\title{Supplementary Material\\[4pt]
\large \textit{Affordable Real Style Transfer: Training Wavelet-Decoupled Rectified Flow on a Single RTX 3060 in Minutes}}
\author{Anonymous Submission}
\affiliations{}
\begin{document}
\maketitle

\section{KL Divergence Verification of Gaussian Assumption}
\label{sec:appendix-kl}

We verify the Gaussian approximation for WCT by computing the KL divergence between the empirical distribution of each wavelet subband and a fitted multivariate Gaussian, on the Distinct5-WikiArt test set (750 latent tensors).

For each subband $h_i \in \mathbb{R}^{4 \times 16 \times 16}$, we flatten the spatial dimensions to obtain feature vectors of dimension $d=4$ (the channel dimension). We then fit a $d$-dimensional Gaussian $\mathcal{N}(\hat\mu_i, \hat\Sigma_i)$ and compute the KL divergence
\begin{equation}
D_{\text{KL}}(p_{\text{emp}} \| \mathcal{N}(\hat\mu_i, \hat\Sigma_i))
= \mathbb{E}_{p_{\text{emp}}}[\log p_{\text{emp}} - \log \mathcal{N}(x;\hat\mu_i,\hat\Sigma_i)].
\end{equation}
The results are:

\begin{center}
\begin{tabular}{lcccc}
\toprule
Subband & LL & LH & HL & HH \\
\midrule
KL divergence (nats) & 0.031 & 0.042 & 0.038 & 0.047 \\
\% of samples within $2\sigma$ & 96.2 & 94.8 & 95.3 & 94.1 \\
\bottomrule
\end{tabular}
\end{center}

All subbands have KL divergence below $0.05$ nats, and over 94\% of samples fall within the $2\sigma$ ellipse of the fitted Gaussian. This confirms that the VAE KL regularization successfully Gaussianizes the latent space, making second-order moment matching (WCT) well-justified.

\section{Proof of Theorem 1 (Local Collapse in Euclidean Latents)}
\label{sec:appendix-thm1}

We restate Theorem~1 in the main text and provide a self-contained proof.

\begin{theorem}[Local Collapse]
Let $\theta_0 \in \mathcal{M}_0$ be stationary in the normal directions:
\begin{equation}
\nabla_\perp \mathcal{L}_{\mathrm{FM}}(\theta_0)=\nabla_\perp \mathcal{L}_{\mathrm{sty}}(\theta_0)=0.
\end{equation}
Let $L_f^{\mathrm{low}}$ and $L_f^{\mathrm{high}}$ be the Lipschitz constants of $\nabla\mathcal{L}_{\mathrm{FM}}$ restricted to the low- and high-frequency latent subspaces, and let $L_s$ be the Lipschitz constant of $\nabla\mathcal{L}_{\mathrm{sty}}$.
If $w_f L_f^{\mathrm{high}} > w_s L_s$, then $\theta_0$ is a strict local minimum of $\mathcal{L}_{\mathrm{euc}} = w_f \mathcal{L}_{\mathrm{FM}} + w_s \mathcal{L}_{\mathrm{sty}}$ on the normal space $T^{\perp}_{\theta_0}\mathcal{M}_0$.
\end{theorem}

\begin{proof}
Consider a normal perturbation $\delta \in T^{\perp}_{\theta_0}\mathcal{M}_0$ with $\|\delta\|$ sufficiently small. The stationarity condition implies $\nabla\mathcal{L}_{\mathrm{euc}}(\theta_0)=0$. A Taylor expansion to second order yields
\begin{equation}
\mathcal{L}_{\mathrm{euc}}(\theta_0+\delta) = \mathcal{L}_{\mathrm{euc}}(\theta_0) + \frac{1}{2}\delta^\top \nabla^2_\perp \mathcal{L}_{\mathrm{euc}}(\theta_0)\delta + o(\|\delta\|^2).
\end{equation}
The normal Hessian decomposes by the Lipschitz bounds:
\begin{equation}
\delta^\top \nabla^2_\perp \mathcal{L}_{\mathrm{euc}}(\theta_0)\delta \ge w_f L_f^{\mathrm{high}} \|\delta\|^2 - w_s L_s \|\delta\|^2.
\end{equation}
Here the first term is the lower bound from the flow-matching Hessian restricted to the high-frequency subspace (where style-conditioned components live), and the second term is the upper bound from the style loss. The $1/f$ spectral structure of VAE latents implies $L_f^{\mathrm{low}} \gg L_f^{\mathrm{high}}$, but the collapse condition only requires the inequality in the high-frequency subspace.

By hypothesis, $w_f L_f^{\mathrm{high}} > w_s L_s$, so the quadratic form is strictly positive definite. Hence $\mathcal{L}_{\mathrm{euc}}(\theta_0+\delta) > \mathcal{L}_{\mathrm{euc}}(\theta_0)$ for any sufficiently small $\delta \neq 0$, confirming $\theta_0$ as a strict local minimum.
\end{proof}

\textbf{Remark.} The condition $w_f L_f^{\mathrm{high}} > w_s L_s$ is satisfied by standard CNN backbones because: (1)~the $1/f$ power spectrum of natural images makes $L_f^{\mathrm{high}}$ non-negligible even though $L_f^{\mathrm{low}}$ dominates; (2)~the style-matching term $\mathcal{L}_{\mathrm{sty}}$ is bounded by the CLIP embedding space, which has a fixed Lipschitz constant; (3)~the flow-matching weight $w_f$ is typically larger than $w_s$ in standard training configurations.

\section{Proof of Proposition 1 (Exponential Content Decay)}
\label{sec:appendix-eota}

We restate Proposition~1 in the main text.

\begin{proposition}[Exponential Content Decay]
Let a style injection of strength $\alpha \in (0,1)$ be applied at every solver step for $n$ steps. The residual content fraction after $n$ steps is
\begin{equation}
r_n(\alpha) = (1-\alpha)^n.
\end{equation}
For $n=8$, $r_8(0.5) = 3.9 \times 10^{-3}$. Endpoint-only injection ($n=1$) yields $r_1(\alpha) = 1-\alpha$, which is affine and preserves blend identifiability.
\end{proposition}

\begin{proof}
At each step, the injection operator $I_\alpha$ blends the current latent $z_t$ toward the style reference by a factor $\alpha$:
\begin{equation}
z_t \leftarrow (1-\alpha) z_t + \alpha \cdot \text{style\_signal}.
\end{equation}
The content component of $z_t$ is multiplied by $1-\alpha$ at each step. After $n$ independent injections, the residual content fraction is the product of $n$ such factors:
\begin{equation}
r_n(\alpha) = \prod_{k=1}^{n} (1-\alpha) = (1-\alpha)^n.
\end{equation}
For $n=1$, the residual is $1-\alpha$, which is an affine function of $\alpha$. For $n \gg 1$, the decay is exponential, compressing the operational blending space and making the style-strength/content-preservation tradeoff non-identifiable.
\end{proof}

\section{Subband Energy Distribution on Distinct5-WikiArt}
\label{sec:appendix-spectrum}

We measured the per-subband energy distribution on the Distinct5-WikiArt test set (750 latent tensors, each $4 \times 32 \times 32$).

\begin{center}
\begin{tabular}{lcccc}
\toprule
& LL & LH & HL & HH \\
\midrule
Mean energy fraction (\%) & 91.7 & 3.2 & 3.0 & 2.1 \\
Std of energy fraction (\%) & 2.1 & 0.8 & 0.7 & 0.6 \\
Style-induced relative change (\%) & 4.3 & 11.8 & 10.2 & 15.6 \\
\bottomrule
\end{tabular}
\end{center}

The LL subband carries $91.7\%$ of the total Frobenius norm on average, confirming the $1/f$ spectral structure. The style-induced relative change (computed as the ratio of per-subband energy difference between content and target-style latents to the content energy) is 2--4$\times$ larger in LH/HL/HH than in LL, confirming that style differences are concentrated in the high-frequency subbands.

\end{document}